\documentclass[12pt,phd,oneside,final]{waynestatedissert}
%
% To change the dissertation to a Master's Thesis, include a documentclass
% option such as [masters], [ms], [ma], etc.  Also available are [wsudraft, twoside].
% As a reminder, documentclass options are a
% comma-separated list, e.g. \documentclass[ms,twoside]{waynestatedissert}
% 

%
% Everything between the \documentclass and the \begin{document} is
% called the preamble of the document. Everything between the
% \begin{document} and the \end{document} is called the body of the
% document. Define any additional commands you want here (the preamble)
%

% BibTeX from the BiBTeX Documentation
%\def\BibTeX{{\rm B\kern-.05em{\sc i\kern-.025em b}\kern-.08em
%    T\kern-.1667em\lower.7ex\hbox{E}\kern-.125emX}}

%
% It is better to break up the dissertation into multiple files (e.g.,
% one file per chapter, as well as separate files for the abstract,
% acknowledgements, and vita).  These files are brought into the
% document using \include{} statements.  There will be times, however,
% when you don't want to print the ENTIRE dissertation.  You can limit
% what will actually be printed by using the \includeonly{} statment.
% This contains a list of the files you want printed.  Any file NOT
% listed will not be printed.  However, all page numbers, references,
% etc., will be preserved as though all the files were actually
% printed. For example, the line below would result only in chapters 1
% and 3 being printed (if it were uncommented).
%

%  \renewcommand\typesetChapterTitle[1]{\uppercase{#1}}
\renewcommand\typesetChapterTitle[1]{#1}
\usepackage{fancyhdr}
\usepackage{bm}
\usepackage{graphicx}
\usepackage{amsmath}
\usepackage{hyperref}
\usepackage{slashed}
\usepackage{mathrsfs}
\usepackage{epsfig}
\usepackage{soul}
\usepackage{amssymb}
\usepackage{amsfonts}
\usepackage{latexsym}
\usepackage{amsthm}
\usepackage{longtable}
\usepackage{etex}
%\reserveinserts{18}
\usepackage{morefloats}
\usepackage{lipsum}

\newcommand{\tmmathbf}[1]{\ensuremath{\boldsymbol{#1}}}
\newcommand{\tmop}[1]{\ensuremath{\operatorname{#1}}}

\usepackage[dvipsnames]{xcolor}
\newcommand{\Red}[1]{{\color{red}#1}}

%\usepackage[margin=1.0in]{geometry}
%\newcommand{\avg}[1]{\left< #1 \right>} % for average

\begin{document}

%
% First, declare the parts of your title page 
%
\author{\uppercase{AUTHOR NAME}}
\title{\uppercase{Thesis title}}
\authordegrees{B.S.}  % Degrees thus far, not including this one.
\unit{Department}

\advisorname{Advisor Name}
%\member{Professor XXX}


%
% The following creates the title page
%
\maketitle

% Next, EITHER a copyright or BLANK page.
%
%   The following creates a page used to copyright your dissertation
%
%   BACKGROUND: Even without this copyright page, your dissertation will
%               carry a common-law copyright. However, if your
%               dissertation ends up seeing wide distribution, your
%               common-law copyright is at risk of being expunged.
%               Adding this copyright page prevents that from happening.
%
%               There are NO DOWNSIDES to including a copyright page as
%               your document is automatically copyright by law anyway.
%               However, this copyright page is OPTIONAL. If you get rid
%               of it, uncomment the \blankpage that follows it so that
%               there is a blank page here. The graduate school requires
%               a page here that is either blank or carries the
%               copyright.
%
%   IMPORTANT NOTE: The graduate school requires either a copyright page
%                   here or a BLANK PAGE here. If you get rid of the
%                   copyright, uncomment the \blankpage that follows it.
%                   You should NOT have BOTH uncommented.
%

% If you get rid of \disscopyright, restore the \blankpage line after it
\disscopyright
%\blankpage

%
% UPDATED TEXT (2010):
%  The graduate school does not require an external abstract. If this
%  changes, follow the old instructions below.
%
% HISTORICAL TEXT (1996):
%  Uncomment the three lines below to generate the external abstract.  Two
%  copies of this must be turned in to the graduate school.  These lines can
%  be placed pretty much anywhere, since the page numbering should be
%  independent of the rest of the thesis
%

% \begin{externalabstract}
%   \input{abstract}
% \end{externalabstract}



\begin{dedication}
{Dedication}
\end{dedication}

%
% Bring in Acknowledgement and Vita from separate files named ``ack.tex''
% and ``vita.tex''.
%

\include{headers/ack}

%
% Make the Table of Contents and other good stuff
%

\tableofcontents
\listoftables
\listoffigures

\newpage 

\pagenumbering{arabic}

\fancypagestyle{plain}{%
\fancyhf{} % clear all header and footer fields
\fancyhead[C]{\thepage} % except the center
\renewcommand{\headrulewidth}{0pt}
\renewcommand{\footrulewidth}{0pt}}
\pagestyle{plain}


%
% The following is a list of chapters.  Each is brought in from a
% separate file using the \include{} command.
%

\include{ch1.intro}     % tell them what you are going to tell them 
\include{ch2.framework}
\include{ch3.results}    % what is your solution
\include{ch4.conclusion}       % conclusion stuff


%
% If you have appendices in your dissertation, you will need the
% following, else keep it commented. The following appendices are in
% files called ``app1.tex'', and ``app2.tex'', and they
% look just like any chapter.
%

\appendix
\include{appendix}


%
% The all important bibliography file at the end of your document!! Use
% the bibstyle you (your department) like in the \bibliographystyle{}
% statement and list the name of your bibliography database file in
% the \bibliography{} statement.  In this example, ``bibfile.bib'' is
% the name of the database.  See the LaTeX manual appendix B for details
% about the bibliography database and BibTeX.
%

%\bibliographystyle{plain}
%\bibliographystyle{hunsrt}
\bibliographystyle{JHEP}
\bibliography{references}

%
% Abstract goes here.
\newpage
\cleardoublepage\phantomsection\addcontentsline{toc}{chapter}{ABSTRACT}
\include{headers/abstract}

\newpage
\cleardoublepage\phantomsection\addcontentsline{toc}{chapter}{AUTOBIOGRAPHICAL STATEMENT}
\include{headers/vita}

\end{document}




